\documentclass[aspectratio=169]{beamer}
\usetheme{Madrid}
\usecolortheme{seagull}
\setbeamertemplate{navigation symbols}{}
\setbeamertemplate{footline}[frame number]

% ---- Paquetes ----
\usepackage[T1]{fontenc}
\usepackage[utf8]{inputenc}
\usepackage[spanish]{babel}
\usepackage{amsmath,amssymb,graphicx,booktabs,xcolor}

% ---- Mismo estilo de main_presentacion_mds ----
\definecolor{myblue}{HTML}{0D47A1}
\setbeamercolor{background canvas}{bg=white}
\setbeamercolor{structure}{fg=myblue}
\setbeamercolor{frametitle}{fg=white,bg=myblue}
\setbeamercolor{headline}{bg=myblue}
\setbeamertemplate{headline}{%
  \begin{beamercolorbox}[wd=\paperwidth,ht=1.8ex,dp=0.4ex]{headline}%
  \end{beamercolorbox}%
}
\setbeamertemplate{frametitle}{%
  \nointerlineskip%
  \begin{beamercolorbox}[wd=\paperwidth,ht=2.6ex,dp=1.2ex,left]{frametitle}
    \hspace*{1ex}\insertframetitle
  \end{beamercolorbox}
}

\graphicspath{{../attachments/Analisis de Datos/figuras_tesis/}}

% ---- Datos ----
\title[Avances Tesis]{Presentación de Avances}
\subtitle{Sincronización y EDA de señales multimodales}
\author[Mario Romero]{Mario Romero}
\institute{Magíster en Data Science}
\date{Agosto 2026}

\begin{document}

\begin{frame}
  \titlepage
\end{frame}

\begin{frame}{Contenidos}
\begin{enumerate}
  \item Contexto y objetivo
  \item Pipeline implementado
  \item Coordinación y sincronización
  \item Resultados del EDA multimodal
  \item Limitaciones metodológicas
  \item Próximos pasos
\end{enumerate}
\end{frame}

% ===== CONTEXTO =====
\section{Contexto}
\begin{frame}{Contexto y objetivo de la tesis}
\textbf{Problema:} atención y activación emocional cambian durante la navegación web;
resumir las señales mediante promedios elimina esa dinámica temporal.

\medskip
\begin{block}{Objetivo}
Comparar modelos temporales multimodales con modelos estáticos equivalentes para predecir
\textbf{atención} y \textbf{activación} durante la navegación web.
\end{block}

\begin{itemize}
  \item Señales disponibles: \textbf{EEG, GSR, pupilometría y eye tracking}.
  \item La Toma de Muestra 2 \textbf{no incluye expresión facial}.
  \item Antes de modelar se debe establecer qué sesiones son utilizables y cómo se coordinan temporalmente.
\end{itemize}
\end{frame}

% ===== PIPELINE =====
\section{Pipeline}
\begin{frame}{Pipeline reproducible implementado}
\begin{columns}[t]
\column{0.48\textwidth}
\begin{enumerate}
  \item \textbf{Inventario} de sesiones EEG y grabaciones Tobii.
  \item \textbf{Normalización} de participantes y tiempos.
  \item \textbf{Coordinación} de sesiones mediante UTC.
\end{enumerate}
\column{0.48\textwidth}
\begin{enumerate}
  \setcounter{enumi}{3}
  \item \textbf{EDA} con criterios específicos por modalidad.
  \item Generación de \textbf{manifiestos, diagnósticos y tablas}.
  \item Exportación de \textbf{seis figuras} reproducibles.
\end{enumerate}
\end{columns}

\begin{alertblock}{Principio de trabajo}
Los archivos originales no se modifican. La lectura se realiza por streaming y las decisiones
quedan documentadas.
\end{alertblock}
\end{frame}

% ===== SINCRONIZACION =====
\section{Sincronización}
\begin{frame}{Coordinación EEG--Tobii}
\begin{columns}[t]
\column{0.42\textwidth}
\begin{block}{Resultado}
\centering
{\Huge \textbf{47}}\\
candidatos directos EEG--Tobii
\end{block}

\begin{itemize}
  \item P29 requiere revisión por offset.
  \item \texttt{Test\_Participant} no tiene pareja directa.
\end{itemize}

\column{0.54\textwidth}
\begin{block}{Ventana común}
\begin{align*}
t_{inicio} &= \max(t_{inicio}^{EEG},t_{inicio}^{Tobii}) \\
t_{fin} &= \min(t_{fin}^{EEG},t_{fin}^{Tobii})
\end{align*}
\end{block}
Cada señal conserva su frecuencia nativa; no es necesario remuestrear para coordinar.
\end{columns}
\end{frame}

\begin{frame}{Límite de la sincronización}
\begin{alertblock}{No existen triggers compartidos}
La alineación absoluta puede estimarse mediante los relojes del sistema, pero el \textbf{drift}
entre dispositivos no puede identificarse válidamente sin eventos observados por ambos equipos.
\end{alertblock}

\begin{itemize}
  \item \texttt{drift\_scale = 1.0} significa ``sin corrección identificable''.
  \item No demuestra ausencia de drift.
  \item Una corrección afín requiere al menos dos anclajes compartidos: TTL, LSL o artefactos deliberados.
  \item Toda fila distinta de \texttt{candidato} debe revisarse manualmente.
\end{itemize}
\end{frame}

% ===== EDA GENERAL =====
\section{EDA multimodal}
\begin{frame}{Disponibilidad y calidad de las señales}
\begin{center}
\includegraphics[width=0.93\linewidth,height=0.78\textheight,keepaspectratio]
{01_resumen_calidad_modalidades.png}
\end{center}
\end{frame}

% ===== EEG =====
\begin{frame}{EDA EEG: 45 de 48 sesiones utilizables}
\begin{columns}
\column{0.57\textwidth}
\includegraphics[width=\linewidth,height=0.67\textheight,keepaspectratio]
{03_eeg_resumen_calidad.png}
\column{0.40\textwidth}
\begin{itemize}
  \item 104 sesiones inventariadas.
  \item Resumen final sobre 48 sesiones principales.
  \item Frecuencia efectiva estable alrededor de \textbf{125 Hz}.
  \item Se excluyen P1, P14 y P17.
\end{itemize}
\begin{alertblock}{Interpretación}
El umbral de 30\,\% es exploratorio y permisivo. ``Utilizable'' no equivale a señal limpia.
\end{alertblock}
\end{columns}
\end{frame}

% ===== GSR =====
\begin{frame}{EDA GSR: frecuencia efectiva cercana a 15 Hz}
\begin{columns}
\column{0.55\textwidth}
\begin{itemize}
  \item 2.984.636 filas procesadas por streaming.
  \item 530.407 muestras GSR válidas.
  \item 47 de 49 grabaciones disponibles: \textbf{95,92\,\%}.
  \item P7 y P42 no contienen GSR.
\end{itemize}
\begin{block}{Corrección respecto del supuesto inicial}
La frecuencia empírica es cercana a \textbf{15 Hz}, no 8 Hz.
\end{block}
\column{0.42\textwidth}
\includegraphics[width=\linewidth,height=0.65\textheight,keepaspectratio]
{04_gsr_resumen.png}
\end{columns}
\end{frame}

% ===== OCULAR =====
\begin{frame}{EDA de pupilometría y eye tracking}
\begin{columns}
\column{0.55\textwidth}
\includegraphics[width=\linewidth,height=0.66\textheight,keepaspectratio]
{05_eye_pupil_validez_participantes.png}
\column{0.42\textwidth}
\begin{itemize}
  \item 2.187.110 muestras oculares evaluadas.
  \item 1.973.941 con ambos ojos válidos: \textbf{90,25\,\%}.
  \item 2.077.373 muestras de mirada válidas: \textbf{94,98\,\%}.
  \item 122.302 fijaciones únicas.
  \item Duración mediana: \textbf{183 ms}.
\end{itemize}
\end{columns}
\end{frame}

% ===== LECTURA CONJUNTA =====
\begin{frame}{Lectura conjunta del avance}
\begin{block}{Resultado principal}
Las cuatro modalidades presentan más de 90\,\% de disponibilidad o validez global, lo que
permite avanzar hacia la construcción de ventanas multimodales.
\end{block}

\begin{alertblock}{Lo que aún no está resuelto}
\begin{itemize}
  \item La calidad global no reemplaza criterios por participante, sesión y canal.
  \item Los emparejamientos no confirmados requieren revisión manual.
  \item La sincronización submuestra no puede defenderse sin triggers compartidos.
  \item La latencia de SCR debe incorporarse al diseño temporal de GSR.
\end{itemize}
\end{alertblock}
\end{frame}

% ===== PROXIMOS PASOS =====
\section{Siguientes pasos}
\begin{frame}{Siguientes pasos}
\begin{enumerate}
  \item Aprobar o corregir el manifiesto de sincronización.
  \item Definir criterios finales de calidad y exclusión.
  \item Recortar ventanas comunes por evento.
  \item Preprocesar cada modalidad conservando su frecuencia nativa.
  \item Construir \textit{Attention Score} y \textit{Arousal Score}.
  \item Entrenar los primeros modelos baseline.
\end{enumerate}

\begin{block}{Decisión inmediata}
Cerrar emparejamientos y criterios de exclusión antes de extraer características o entrenar modelos.
\end{block}
\end{frame}

\begin{frame}{}
\centering
\vspace{1cm}
{\Large \textbf{Gracias}}\\[1em]
¿Preguntas?
\end{frame}

\end{document}
